\begin{lemma}[Invariance of the five predicates under modification on a null set]
Let $n\ge1$ be an integer and $q\ge1$ real, let $\Omega\subseteq\R^n$ be open,
let $f,f'\colon\R^n\to\R$ and $g\colon\R^n\to\R^n$, and suppose
\[
f=f'\qquad\text{almost everywhere on }\Omega .
\]
Then each of the following five implications holds, with the \emph{same} field
$g$ on both sides:
\begin{enumerate}
\item if $g$ is the weak gradient of $f$ on $\Omega$
  (\noderef{HasWeakGradientOn}), then $g$ is the weak gradient of $f'$ on
  $\Omega$;
\item if $f\in W^{1,q}(\Omega)$ with weak gradient $g$ (\noderef{MemW1pOn}),
  then $f'\in W^{1,q}(\Omega)$ with weak gradient $g$;
\item if $f\in W^{1,q}_{\mathrm{loc}}(\Omega)$ with weak gradient $g$
  (\noderef{MemW1pLocOn}), then $f'\in W^{1,q}_{\mathrm{loc}}(\Omega)$ with weak
  gradient $g$;
\item for all $A\colon\R^n\times\R\times\R^n\to\R^n$ and
  $B\colon\R^n\times\R\times\R^n\to\R$, if $f$ is a weak solution of
  (1.1) of the thesis in $\Omega$ at the exponent $q$ with weak gradient $g$
  (\noderef{IsWeakSolutionOn}), then so is $f'$; and
\item for every $f^{\sharp}\colon\R^n\times\R\times\R^n\to\R$, if $f$ is a local
  minimum of $F(\,\cdot\,;E)=\int_Ef^{\sharp}(x,\cdot,\nabla\cdot)\,dx$ in
  $\Omega$ at the exponent $q$ with weak gradient $g$
  (\noderef{IsLocalMinimumOn}), then so is $f'$.
\end{enumerate}
Here ``at the exponent $q$'' refers to the test class and the local
integrability that \noderef{IsWeakSolutionOn} and \noderef{IsLocalMinimumOn}
carry, namely $W^{1,q}_{\mathrm{loc}}$ membership of the function and
$W^{1,q}(\Omega)$ test pairs of compact support.

The five clauses are one node because they are one fact used in one way.  In
each of the five predicates the function enters \emph{only} through integrals
over $\Omega$, or over a compact subset of $\Omega$, or over the support of a
test function contained in $\Omega$; the integrand is changed only on the null
set $\Omega\cap\{f\ne f'\}$, and an integral does not see a null set.  No clause
is a special case of another: (1) is about the identity alone, (2) and (3) add
the two different integrability packages, and (4) and (5) are about the two
different nonlinear functionals of the function's \emph{values}, where the
substitution happens inside $A$, $B$ and $f^{\sharp}$ rather than in a factor.

\medskip
\emph{What it is for.}  The paper replaces $u$ by its continuous representative
without comment: having obtained that ``$u$ is locally H\"older continuous''
(p.~31) it writes $f^0(h)=f(x_0,u(x_0),h)$ (p.~32) and treats $u$ as defined
pointwise.  Formally a continuous representative obtained from an oscillation
estimate agrees with $u$ only almost everywhere, so
every hypothesis carried on $u$ --- the Sobolev membership, the equation, the
minimum property --- has to be transported to $U$ before the frozen operator and
the comparison map can be formed.  That transport is this lemma.  Nothing here
is specific to the operators of the paper: the statement mentions no structure
condition and no estimate.
\end{lemma}
\begin{proof}
Throughout, ``test function'' means a $\varphi\in C^\infty(\R^n)$ with compact
support and $\operatorname{supp}\varphi\subseteq\Omega$, which is the test class
of \noderef{HasWeakGradientOn}, and $N=\Omega\cap\{x:f(x)\ne f'(x)\}$, which is
a Lebesgue null set: the hypothesis says exactly that $N$ is null for Lebesgue
measure restricted to $\Omega$, and that measure agrees with Lebesgue measure on
measurable subsets of $\Omega$.

We use two standard facts about the Bochner integral, both of which hold because
the integral depends only on the almost-everywhere class of the integrand with
respect to the measure in question, and neither of which needs an integrability
hypothesis:
\begin{itemize}
\item[(A1)] if $v=v'$ almost everywhere with respect to a measure $\mu$, then
  $\int v\,d\mu=\int v'\,d\mu$ (in Mathlib,
  \texttt{MeasureTheory.integral\_congr\_ae}; when neither integral exists both
  are $0$ by the convention for a non-integrable function, so the identity holds
  in every case);
\item[(A2)] if $v=v'$ almost everywhere with respect to $\mu$ and
  $v\in L^{q}(\mu)$, then $v'\in L^{q}(\mu)$, with the same $L^q$ norm
  (in Mathlib, \texttt{MeasureTheory.MemLp.ae\_eq}; almost-everywhere strong
  measurability and the finiteness of the norm are both properties of the
  almost-everywhere class).
\end{itemize}
We also use that almost-everywhere statements pass to measurable subsets: if
$E'\subseteq E$ are measurable and a property holds almost everywhere on $E$,
it holds almost everywhere on $E'$, because the measure restricted to $E'$ is
dominated by the measure restricted to $E$ (Preamble, the transport facts
recorded with H\"older's inequality, \texttt{Measure.restrict\_mono}).

\emph{Step 1: clause (1).}  Assume $\mathrm{HasWeakGradientOn}(n,\Omega,f,g)$,
let $\varphi$ be a test function and let $i\in\{1,\dots,n\}$.  The two functions
$x\mapsto f(x)D_i\varphi(x)$ and $x\mapsto f'(x)D_i\varphi(x)$ agree at every
point of $\Omega\setminus N$, hence almost everywhere on $\Omega$, so (A1) for
Lebesgue measure restricted to $\Omega$ gives
\[
\int_\Omega f'\,D_i\varphi\,dx=\int_\Omega f\,D_i\varphi\,dx
=-\int_\Omega g^i\,\varphi\,dx ,
\]
the second equality being the hypothesis.  As $\varphi$ and $i$ were arbitrary,
this is $\mathrm{HasWeakGradientOn}(n,\Omega,f',g)$.

\emph{Step 2: clause (2).}  By \noderef{MemW1pOn} the hypothesis is the
conjunction of $\mathrm{HasWeakGradientOn}(n,\Omega,f,g)$, of
$f\in L^q(\Omega)$ and of $|g|\in L^q(\Omega)$.  Step 1 upgrades the first to
$f'$; (A2), applied to Lebesgue measure restricted to $\Omega$, upgrades the
second to $f'\in L^q(\Omega)$; and the third does not mention $f$.  By
\noderef{MemW1pOn} again, $f'\in W^{1,q}(\Omega)$ with weak gradient $g$.

\emph{Step 3: clause (3).}  By \noderef{MemW1pLocOn} the hypothesis is the
conjunction of $\mathrm{HasWeakGradientOn}(n,\Omega,f,g)$ with: for every
compact $K\subseteq\Omega$, $f\in L^q(K)$ and $|g|\in L^q(K)$.  Step 1 upgrades
the identity.  For the local clause, fix a compact $K\subseteq\Omega$; then
$f=f'$ almost everywhere on $K$, because $K\subseteq\Omega$ and
almost-everywhere statements pass to measurable subsets, so (A2) gives
$f'\in L^q(K)$, while $|g|\in L^q(K)$ is unchanged.  By \noderef{MemW1pLocOn},
$f'\in W^{1,q}_{\mathrm{loc}}(\Omega)$ with weak gradient $g$.

\emph{Step 4: clause (4).}  Fix $A$ and $B$ and assume $f$ is a weak solution
with weak gradient $g$.  By \noderef{IsWeakSolutionOn} this is the conjunction of
$f\in W^{1,q}_{\mathrm{loc}}(\Omega)$ with weak gradient $g$ --- upgraded to
$f'$ by Step 3 --- and, for each admissible test pair, of the integrability on
$\Omega$ of the integrand below together with the identity
\[
\int_\Omega\Bigl(\sum_{i=1}^n\bigl(A(x,f(x),g(x))\bigr)^i\,(g_\varphi(x))^i
+B(x,f(x),g(x))\,\varphi(x)\Bigr)dx=0
\]
for every pair $(\varphi,g_\varphi)$ with $\varphi\in W^{1,q}(\Omega)$ of weak
gradient $g_\varphi$ (\noderef{MemW1pOn}), $\varphi$ of compact support and
$\operatorname{supp}\varphi\subseteq\Omega$.  Fix such a pair.  At every point
of $\Omega\setminus N$ we have $f(x)=f'(x)$, so the integrand formed with $f'$
equals the integrand formed with $f$ there --- the substitution occurs only in
the second argument of $A$ and of $B$, and the remaining arguments $x$, $g(x)$,
$g_\varphi(x)$, $\varphi(x)$ do not involve $f$ --- hence the two integrands
agree almost everywhere on $\Omega$.  Both clauses transfer on that one
observation: integrability on $\Omega$ is preserved under altering a function on
a null set, so the $f'$-integrand is integrable because the $f$-integrand is,
and (A1) gives the same integral, namely $0$.  As the pair was arbitrary, $f'$
is a weak solution with weak gradient $g$ by \noderef{IsWeakSolutionOn}.

\emph{Step 5: clause (5).}  Fix $f^{\sharp}$ and assume $f$ is a local minimum
with weak gradient $g$.  By \noderef{IsLocalMinimumOn} this is the conjunction of
$f\in W^{1,q}_{\mathrm{loc}}(\Omega)$ with weak gradient $g$ --- upgraded to
$f'$ by Step 3 --- and, for each admissible $\varphi$, of the integrability on
$\operatorname{supp}\varphi$ of \emph{both} integrands below together with the
inequality
\[
\int_{\operatorname{supp}\varphi}f^{\sharp}\bigl(x,f(x),g(x)\bigr)dx
\ \le\
\int_{\operatorname{supp}\varphi}f^{\sharp}\bigl(x,f(x)+\varphi(x),g(x)+g_\varphi(x)\bigr)dx
\]
for every admissible pair $(\varphi,g_\varphi)$ as in Step 4.  Fix such a pair
and put $E=\operatorname{supp}\varphi$, a measurable subset of $\Omega$.  Since
$f=f'$ almost everywhere on $\Omega$, the same holds almost everywhere on $E$.
Hence the integrand on the left formed with $f'$ agrees almost everywhere on $E$
with the one formed with $f$, and likewise on the right, where the substitution
occurs in $f(x)+\varphi(x)$ only.  Both clauses transfer on that one
observation: integrability on $E$ is preserved under altering a function on a
null subset of $E$, so each of the two $f'$-integrands is integrable on $E$
because the corresponding $f$-integrand is; and (A1), applied to Lebesgue
measure restricted to $E$, leaves both sides unchanged, so the inequality holds
for $f'$.  As the pair was arbitrary, $f'$ is a local minimum with weak gradient $g$
by \noderef{IsLocalMinimumOn}.
\end{proof}
